\documentclass[11pt,a4paper]{article}

\usepackage{times}
\usepackage{latexsym}
\usepackage{amsmath}
\usepackage{amssymb}
\usepackage{booktabs}
\usepackage{multirow}
\usepackage{graphicx}
\usepackage{hyperref}
\usepackage{url}
\usepackage{microtype}
\usepackage{xcolor}
\usepackage{enumitem}
\usepackage{algorithm}
\usepackage{algorithmic}
\usepackage{geometry}
\usepackage{array}
\usepackage{tabularx}
\usepackage{caption}
\usepackage{subcaption}
\usepackage{listings}
\usepackage{setspace}

\geometry{margin=1in}
\hypersetup{colorlinks=true, linkcolor=blue, citecolor=blue, urlcolor=blue}

\lstset{
  basicstyle=\small\ttfamily,
  breaklines=true,
  frame=single,
  backgroundcolor=\color{gray!10},
}

\definecolor{primary}{RGB}{31,119,180}
\definecolor{supporting}{RGB}{44,160,44}
\definecolor{contrasting}{RGB}{214,39,40}
\definecolor{qualifying}{RGB}{255,127,14}
\definecolor{superseded}{RGB}{148,103,189}

% ── Title ─────────────────────────────────────────────────────────────────────

\title{
  \textbf{PRISM: Epistemic Graph Retrieval via Typed Spreading Activation} \\[0.5em]
  \large Structured Knowledge Retrieval Through Convergence Scoring and Epistemic Bucketing
}

\author{
  \textbf{MadMando}\thanks{\href{https://github.com/MadMando}{github.com/MadMando}} \\
  \textbf{Claude (Anthropic)}\thanks{Claude Sonnet 4.6 / Opus 4.7, contributed to the adapter architecture, the multi-backend refactor, and this manuscript.} \\[0.5em]
  \href{https://github.com/MadMando/prism}{https://github.com/MadMando/prism}
}

\date{April 2026 — revised for v0.2.5}

\begin{document}

\maketitle

% ── Abstract ──────────────────────────────────────────────────────────────────

\begin{abstract}
Retrieval-Augmented Generation (RAG) systems conventionally retrieve a flat ranked
list of semantically similar passages, treating all retrieved knowledge as
epistemically equivalent. This ignores the rich relational fabric of real corpora:
some passages \emph{support} a claim, others \emph{refute} it, still others
\emph{qualify} or \emph{supersede} it. We introduce \textbf{PRISM}
(Propagation \& Retrieval via Informed Semantic Mapping), a retrieval framework
that layers a \emph{typed epistemic knowledge graph} over an existing vector store
and uses \emph{spreading activation} with \emph{convergence scoring} to retrieve
knowledge structured by how it relates, not merely by how similar it is.

PRISM defines a ten-type epistemic edge taxonomy
(\textsc{supports}, \textsc{refutes}, \textsc{supersedes}, \textsc{derives\_from},
\textsc{specializes}, \textsc{contrasts\_with}, \textsc{implements},
\textsc{generalizes}, \textsc{exemplifies}, \textsc{qualifies}),
each assigned a propagation weight and an epistemic valence.
At retrieval time, a query activates seed nodes via vector similarity; activation
then propagates through the typed graph for $h$ hops with per-type decay.
Nodes reached by \emph{multiple independent seed paths} receive a convergence
bonus — operationalising the memory consolidation insight that knowledge
accessible from many contexts is more reliably grounded. Results are returned
not as a ranked list but as five structured \emph{epistemic buckets}:
\textsc{primary}, \textsc{supporting}, \textsc{contrasting},
\textsc{qualifying}, and \textsc{superseded}.

PRISM requires no re-indexing of existing vector stores, falls back to pure
vector retrieval when no graph is available, and ships as an open-source
Python library (\texttt{pip install prism-rag}).
We describe the architecture, the edge taxonomy, the spreading activation
algorithm, and preliminary evaluation on a synthetic multi-source benchmark.
\end{abstract}

% ── 1. Introduction ───────────────────────────────────────────────────────────

\section{Introduction}
\label{sec:intro}

Retrieval-Augmented Generation \citep{lewis2020rag} has become the dominant
paradigm for grounding large language model (LLM) outputs in external knowledge.
A standard RAG pipeline embeds a query, retrieves the top-$K$ most similar
passages from a vector store, and concatenates them as context for the generating
model. The implicit assumption is that passages can be ranked by a single scalar
similarity score — that relevance is ordinal.

This assumption breaks down in corpora that contain \emph{conflicting},
\emph{temporally stratified}, or \emph{structurally related} knowledge.
Consider a regulatory document corpus spanning ten years: the most recent
standard may \emph{supersede} an older one on the same topic; a practitioner
guide may \emph{qualify} a general principle with important exceptions; two
competing frameworks may both be highly similar to a query while
\emph{refuting} each other. A flat ranked list delivers all of this to the LLM
without structural annotation. The LLM must then implicitly reason about
relationships that the retrieval layer discarded.

\textbf{Our position:} Retrieval is not a ranking problem. It is a
\emph{knowledge structuring} problem. The correct output of a retrieval step
is not \enquote{the top-5 most similar passages} but rather
\enquote{here is what supports the answer, here is what challenges it,
here is what qualifies it, here is what has been superseded.}

We introduce \textbf{PRISM}, a retrieval framework that operationalises
this position through three components:

\begin{enumerate}[leftmargin=*]
  \item \textbf{Epistemic Graph Construction.} An offline phase in which
        an LLM classifies relationships between semantically proximate
        passage pairs, producing a directed multigraph with typed, weighted edges.

  \item \textbf{Typed Spreading Activation.} A query-time activation
        algorithm inspired by Collins \& Loftus's spreading activation theory
        \citep{collins1975spreading} in which activation propagates from
        vector-search seed nodes through typed edges, with per-type propagation
        weights that reflect the epistemic strength of each relationship.

  \item \textbf{Convergence Scoring \& Epistemic Bucketing.} Nodes reached
        by multiple \emph{independent} seed paths receive a convergence bonus;
        nodes are then assigned to one of five epistemic buckets based on the
        dominant valence of the edges through which they were reached.
\end{enumerate}

PRISM makes no assumptions about the internal structure of the underlying
vector store. Built-in adapters ship for LanceDB, ChromaDB, Qdrant,
Weaviate, and pgvector; a documented \texttt{VectorAdapter} protocol
lets users plug in any other backend without re-embedding or re-indexing
their corpus. When no graph is available, PRISM degrades gracefully to
standard vector retrieval.

\paragraph{Contributions.}
\begin{itemize}[leftmargin=*]
  \item A ten-type \emph{epistemic edge taxonomy} with formal definitions,
        propagation weights, and valence assignments (\S\ref{sec:taxonomy}).
  \item A \emph{convergence-scoring} variant of spreading activation that
        rewards nodes accessible from multiple independent query perspectives
        (\S\ref{sec:activation}).
  \item An \emph{epistemic bucketing} output schema that structures
        retrieval results for downstream LLM consumption (\S\ref{sec:bucketing}).
  \item An open-source Python library (\texttt{prism-rag} on PyPI, v0.2.7)
        with a pluggable \texttt{VectorAdapter} protocol, built-in adapters
        for LanceDB, ChromaDB, Qdrant, Weaviate, and pgvector, diagnostic
        CLIs (\texttt{prism-build}, \texttt{prism-stats}, \texttt{prism-inspect},
        \texttt{prism-viz}, \texttt{prism-export}), NetworkX and Neo4j export,
        a local interactive graph explorer (\texttt{prism-explore}),
        and a self-contained synthetic benchmark.
\end{itemize}

% ── 2. Background ─────────────────────────────────────────────────────────────

\section{Background and Related Work}
\label{sec:background}

\subsection{Retrieval-Augmented Generation}

RAG \citep{lewis2020rag} augments an LLM with a retrieval component,
conditioning generation on retrieved passages rather than relying solely on
parametric knowledge. Dense retrieval using bi-encoder models
\citep{karpukhin2020dpr} has largely supplanted sparse BM25 retrieval for
semantic similarity. Iterative retrieval methods such as IRCoT
\citep{trivedi2022ircot} interleave retrieval and reasoning, improving
multi-hop question answering at the cost of latency. PRISM targets single-step
retrieval but provides richer structural context than flat dense retrieval.

\subsection{Graph-Augmented RAG}

Microsoft GraphRAG \citep{edge2024graphrag} extracts entity and community
structures from a corpus and uses them to support both local (entity-centric)
and global (summary-centric) query modes. LightRAG \citep{guo2024lightrag}
improves efficiency with a dual-level retrieval architecture over
entity-relationship graphs. Neither system assigns \emph{epistemic types}
to edges; relationships are treated as generic \enquote{relates to} links with
no propagation semantics.

\subsection{HippoRAG}

The most architecturally proximate prior work is HippoRAG
\citep{gutierrez2024hipporag}, which draws on the hippocampal indexing theory
of human long-term memory \citep{teyler1986hippocampal} to motivate a
retrieval architecture that combines LLM-extracted knowledge graphs with
Personalized PageRank (PPR). HippoRAG demonstrates substantial gains on
multi-hop QA benchmarks (up to 20\% over state-of-the-art methods) and
shows that single-step graph-augmented retrieval can match iterative retrieval
at a fraction of the cost.

PRISM shares HippoRAG's neurobiological motivation but differs in three
critical ways. First, HippoRAG's graph is built from \emph{entity co-occurrence}
(named entity pairs extracted from passages); PRISM's graph encodes
\emph{epistemic relationships between passages} with a typed taxonomy.
Second, HippoRAG uses Personalized PageRank, which assigns a global importance
score per node; PRISM uses spreading activation, which is
\emph{source-aware} — it tracks which seed nodes independently activate which
graph nodes, enabling convergence scoring. Third, HippoRAG returns a ranked
list; PRISM returns structured epistemic buckets.

\subsection{Spreading Activation Theory}

Collins \& Loftus \citep{collins1975spreading} proposed spreading activation
as a model of human semantic memory: concepts are nodes in a network,
associations are weighted edges, and activation propagates from cued concepts
to related ones with decay proportional to path length. PRISM directly
instantiates this model with a query-derived initial activation distribution,
per-type propagation weights, and a hop-based decay schedule.

\subsection{Knowledge Representation and Epistemic Logic}

The PRISM edge taxonomy draws on work in epistemic logic
\citep{fagin2003reasoning} and knowledge representation
\citep{brachman1985overview}. The distinction between \emph{positive},
\emph{dialectical}, \emph{qualifying}, and \emph{temporal} edge valences
echoes the defeasible reasoning literature \citep{pollock1995cognitive},
in which arguments can be \emph{undercut} (qualified), \emph{rebutted}
(refuted), or \emph{defeated} (superseded).

% ── 3. The PRISM Framework ────────────────────────────────────────────────────

\section{The PRISM Framework}
\label{sec:framework}

\subsection{Problem Formulation}

Let $\mathcal{C} = \{p_1, p_2, \ldots, p_m\}$ be a corpus of text passages,
each associated with a source document $\text{src}(p_i)$.
Let $E : \mathcal{C} \to \mathbb{R}^d$ be an embedding function.
Given a query $q$, standard dense retrieval returns
$\text{top-}K\{E(q) \cdot E(p_i)\}$.

PRISM augments this with an \emph{epistemic graph}
$G = (\mathcal{C}, \mathcal{E})$ where each edge
$e_{ij} = (p_i, p_j, r, c) \in \mathcal{E}$ carries a
\emph{relation type} $r \in \mathcal{R}$ and a
\emph{confidence score} $c \in [0, 1]$.
Retrieval then returns a structured result
$\mathcal{O} = \{\mathcal{B}_\text{primary},\ \mathcal{B}_\text{supporting},\
\mathcal{B}_\text{contrasting},\ \mathcal{B}_\text{qualifying},\
\mathcal{B}_\text{superseded}\}$.

\subsection{Epistemic Edge Taxonomy}
\label{sec:taxonomy}

We define a ten-type relation taxonomy $\mathcal{R}$, shown in
Table~\ref{tab:taxonomy}. Each type is assigned a
\emph{propagation weight} $w(r) \in (0, 1]$ that controls how strongly
activation flows across edges of that type, and an
\emph{epistemic valence} $v(r) \in \{\text{positive}, \text{qualifying},
\text{dialectical}, \text{temporal}\}$ that determines bucket assignment.

\begin{table}[ht]
\centering
\caption{Epistemic edge taxonomy with propagation weights and valences.}
\label{tab:taxonomy}
\small
\begin{tabular}{llcp{4.0cm}}
\toprule
\textbf{Type} & \textbf{Valence} & \textbf{Weight} & \textbf{Semantics} \\
\midrule
\textsc{supports}       & positive    & 0.90 & $p_i$ provides evidence reinforcing $p_j$ \\
\textsc{derives\_from}  & positive    & 0.80 & $p_i$ is logically derived from $p_j$ \\
\textsc{generalizes}    & positive    & 0.75 & $p_i$ is a broader abstraction of $p_j$ \\
\textsc{specializes}    & positive    & 0.75 & $p_i$ is a specific instance of $p_j$ \\
\textsc{implements}     & positive    & 0.70 & $p_i$ concretely realises the concept in $p_j$ \\
\textsc{exemplifies}    & positive    & 0.65 & $p_i$ is a concrete example illustrating $p_j$ \\
\textsc{qualifies}      & qualifying  & 0.75 & $p_i$ adds conditions or exceptions to $p_j$ \\
\textsc{contrasts\_with}& dialectical & 0.60 & $p_i$ and $p_j$ take different positions \\
\textsc{refutes}        & dialectical & 0.60 & $p_i$ directly contradicts $p_j$ \\
\textsc{supersedes}     & temporal    & 0.40 & $p_i$ replaces $p_j$ (newer / more authoritative) \\
\bottomrule
\end{tabular}
\end{table}

The weight ordering reflects epistemic strength: \textsc{supports}
carries activation most strongly because it represents corroborating evidence;
\textsc{supersedes} carries it most weakly because historical context,
while relevant, is less actionable than current material.

\subsection{Graph Construction}
\label{sec:construction}

Graph construction is an offline phase executed once per corpus (or
incrementally when new sources are added).

\paragraph{Step 1: Candidate pair generation.}
For each passage $p_i$, we retrieve its $k$-nearest neighbours by
cosine similarity over the embedded corpus:
\[
  N(p_i) = \text{top-}k \left\{ \cos(E(p_i), E(p_j)) : j \neq i \right\}
\]
By default, we restrict candidates to \emph{cross-source} pairs
($\text{src}(p_i) \neq \text{src}(p_j)$). This focuses extraction on
inter-document epistemic relationships, which carry the most signal for
structured retrieval. Intra-document relationships tend to reflect
rhetorical continuity rather than epistemic contrast.

\paragraph{Step 2: LLM-based relation extraction.}
Candidate pairs are batched (default batch size: 5) and submitted to an
OpenAI-compatible LLM with a structured prompt:

\begin{lstlisting}
For each pair below, determine if an epistemic relationship
exists. If so, output JSON:
{"pair_idx": N, "type": "<relation_type>",
 "direction": "A->B" | "B->A", "confidence": 0.0-1.0}
Output null for pairs with no clear relationship.
\end{lstlisting}

The LLM returns structured JSON; pairs with confidence $c < \theta$
(default $\theta = 0.65$) are discarded. Validated triples are added as
directed edges to the graph $G$.

\paragraph{Step 3: Graph serialisation.}
The completed graph is serialised to gzip-compressed JSON and stored
alongside the vector store. At retrieval time, the graph is loaded into
memory in a single pass (sub-second for graphs of up to 100k nodes).

\subsection{Typed Spreading Activation}
\label{sec:activation}

Given a query $q$, we initialise a per-node activation score.

\paragraph{Seed initialisation.}
The top-$K_s$ passages by vector similarity are designated seeds:
\[
  a^{(0)}(p_i) = \begin{cases}
    \cos(E(q), E(p_i)) & \text{if } p_i \in \text{seeds} \\
    0 & \text{otherwise}
  \end{cases}
\]

\paragraph{Propagation.}
For $h$ hops, activation propagates from each currently-active node
to its graph neighbours, weighted by edge type and attenuated by a
global decay factor $\delta$:
\[
  a^{(t)}(p_j) = \max_{(p_i, p_j, r, c) \in \mathcal{E}}
    \left[ a^{(t-1)}(p_i) \cdot w(r) \cdot c \cdot \delta^t \right]
\]
We use a max-pool rather than sum-pool to prevent high-degree nodes from
accumulating disproportionate activation from many weak neighbours.

\paragraph{Convergence tracking.}
Throughout propagation, we track the set of seed nodes that have
contributed activation to each node:
\[
  \text{seeds}(p_j) = \left\{ s \in \text{seeds} :
    \exists \text{ path from } s \text{ to } p_j \text{ in } G \right\}
\]
The convergence score is the fraction of seeds that independently
activated node $p_j$:
\[
  \text{conv}(p_j) = \frac{|\text{seeds}(p_j)|}{|\text{seeds}|}
\]

\paragraph{Final scoring.}
The final score combines activation magnitude with convergence:
\[
  f(p_j) = (1 - \lambda) \cdot a(p_j) + \lambda \cdot \text{conv}(p_j)
\]
where $\lambda$ is the convergence weight (default $\lambda = 0.4$).
Nodes with high convergence — reached by multiple independent query
perspectives — are promoted relative to nodes activated strongly by a
single seed.

\subsection{Epistemic Bucketing}
\label{sec:bucketing}

Retrieved nodes are assigned to epistemic buckets based on two criteria:
seed status and the dominant valence of the edges through which they
were reached.

\begin{enumerate}[leftmargin=*]
  \item \textbf{PRIMARY}: Seed nodes (direct vector matches) and
        non-seed nodes with convergence $\geq \tau_\text{conv}$.
        These represent the core, most convergently supported answer.

  \item \textbf{SUPPORTING}: Non-seed nodes reached primarily via
        edges with \emph{positive} valence
        (\textsc{supports}, \textsc{derives\_from}, \textsc{generalizes},
        \textsc{specializes}, \textsc{implements}, \textsc{exemplifies}).

  \item \textbf{CONTRASTING}: Non-seed nodes reached primarily via
        edges with \emph{dialectical} valence
        (\textsc{refutes}, \textsc{contrasts\_with}).

  \item \textbf{QUALIFYING}: Non-seed nodes reached primarily via
        \emph{qualifying} edges (\textsc{qualifies}).

  \item \textbf{SUPERSEDED}: Non-seed nodes reached via
        \emph{temporal} edges (\textsc{supersedes}),
        indicating historically relevant but now-replaced content.
\end{enumerate}

The output is a structured \texttt{EpistemicResult} object containing
all five buckets with per-chunk metadata: source, page, section,
vector score, activation score, convergence score, final score,
and the edge types via which the chunk was reached.

This structure is designed for direct consumption by an LLM prompt:
\begin{lstlisting}
## PRIMARY
[1] source-a  p.14  §2.1  (score: 0.923)
    The core relevant passage...

## CONTRASTING VIEWS
[1] source-b  p.67  §5.0  (score: 0.812  [via: refutes])
    A passage that directly challenges the above...

## QUALIFICATIONS & NUANCES
[1] source-c  p.38  §3.1  (score: 0.712  [via: qualifies])
    Conditions and exceptions to consider...
\end{lstlisting}

By annotating each retrieved passage with its epistemic role, PRISM
provides the LLM with information that flat ranked lists cannot:
\emph{which passages to treat as primary evidence, which to weigh as
counterarguments, and which to treat as historical context only}.

% ── 4. Comparison with Related Systems ───────────────────────────────────────

\section{Comparison with Related Systems}
\label{sec:comparison}

Table~\ref{tab:comparison} summarises how PRISM differs from major
graph-augmented RAG systems.

\begin{table}[ht]
\centering
\caption{Comparison of graph-augmented RAG systems.}
\label{tab:comparison}
\small
\begin{tabular}{lccccc}
\toprule
\textbf{System} & \textbf{Typed} & \textbf{Spreading} & \textbf{Convergence} & \textbf{Epistemic} & \textbf{Builds} \\
                & \textbf{Edges} & \textbf{Activation} & \textbf{Scoring} & \textbf{Buckets} & \textbf{Own Graph} \\
\midrule
Dense RAG            & \ding{55} & \ding{55} & \ding{55} & \ding{55} & \ding{55} \\
GraphRAG             & \ding{55} & \ding{55} & \ding{55} & \ding{55} & \ding{51} \\
LightRAG             & partial   & \ding{55} & \ding{55} & \ding{55} & \ding{51} \\
HippoRAG             & \ding{55} & $\sim$PPR & \ding{55} & \ding{55} & \ding{51} \\
\textbf{PRISM}       & \ding{51} & \ding{51} & \ding{51} & \ding{51} & \ding{51} \\
\bottomrule
\end{tabular}
\end{table}

The key distinctions relative to HippoRAG — the closest prior work —
are as follows.

\paragraph{Edge typing vs. entity co-occurrence.}
HippoRAG builds its graph by extracting named entity mentions and
recording co-occurrence within passages. This captures \emph{topical proximity}
but not \emph{epistemic relationship}. Two passages that both mention
\enquote{data governance} are linked regardless of whether one supports,
refutes, or qualifies the other. PRISM's extraction step explicitly asks
the LLM to classify the epistemic type of each relationship,
producing a graph whose edge types carry semantics that can be used
both for activation weighting and for bucket assignment.

\paragraph{Spreading activation vs. Personalized PageRank.}
PageRank and spreading activation are both diffusion processes on graphs,
but they differ in an important way: PageRank is \emph{global} — it
assigns a single importance score that converges to the stationary
distribution of a random walk, marginalising over all starting points.
Spreading activation is \emph{source-aware} — each seed's contribution
is tracked independently throughout propagation. This source-awareness
is what enables convergence scoring: we can ask not just
\enquote{how activated is this node?} but \enquote{how many
\emph{independent} perspectives on the query converge here?}

\paragraph{Convergence scoring as a reliability signal.}
The convergence score $\text{conv}(p_j)$ measures the fraction of
independent query-relevant seeds that transitively activate node $p_j$.
A passage activated by five independent seeds is more likely to represent
genuinely relevant, cross-validated knowledge than one activated by
a single strong seed path. This is directly analogous to memory
consolidation in hippocampal indexing theory: memories that are cued
from multiple contexts become more stable and accessible
\citep{mcclelland1995complementary}.

% ── 5. Preliminary Evaluation ─────────────────────────────────────────────────

\section{Preliminary Evaluation}
\label{sec:evaluation}

\subsection{Synthetic Benchmark}
\label{sec:synthetic}

To demonstrate epistemic bucketing without requiring a labelled QA dataset,
we construct a synthetic benchmark consisting of 8 passages from 3
fictional sources concerning a shared topic (accountability in distributed
systems). The corpus is designed with known ground-truth epistemic
relationships: one passage \textsc{refutes} another, one \textsc{supersedes}
an earlier treatment, one \textsc{qualifies} a general principle.
A small epistemic graph of 6 edges is constructed directly (without LLM
extraction) and used for retrieval.

\paragraph{Standard vector retrieval result:}
\begin{lstlisting}
Query: "accountability mechanisms in distributed systems"

[1] framework-a: "Accountability requires clear ownership..." (0.912)
[2] framework-b: "Distributed accountability is incompatible..." (0.891)
[3] standard-c:  "Accountability must be assigned at..." (0.856)
[4] framework-a: "Historical note: early frameworks..." (0.821)
[5] framework-b: "An exception applies when..." (0.798)
\end{lstlisting}

\paragraph{PRISM epistemic retrieval result:}
\begin{lstlisting}
Query: "accountability mechanisms in distributed systems"

## PRIMARY
[1] framework-a: "Accountability requires clear ownership..." (0.912)
[2] standard-c:  "Accountability must be assigned at..." (0.856)

## CONTRASTING VIEWS
[1] framework-b: "Distributed accountability is incompatible..."
    (0.891  [via: refutes])

## QUALIFICATIONS & NUANCES
[1] framework-b: "An exception applies when..."
    (0.798  [via: qualifies])

## SUPERSEDED (historical context)
[1] framework-a: "Historical note: early frameworks treated..."
    (0.821  [via: supersedes])
\end{lstlisting}

The epistemic structure transforms a flat list into an answer that
explicitly flags the contradiction (\textsc{refutes}), the conditions
(\textsc{qualifies}), and the historical context (\textsc{supersedes}).
An LLM consuming the second format has the information it needs to
produce a calibrated, nuanced answer; one consuming the first must
discover these relationships implicitly or miss them entirely.

\subsection{Large-Corpus Graph Statistics}
\label{sec:corpus}

We applied PRISM to a proprietary governance knowledge corpus of
30,533 passages from 14 distinct source documents (technical standards,
frameworks, practitioner guides). Using $k=8$ nearest neighbours per
passage with cross-source filtering, we generated 49,571 candidate pairs
and submitted them to an LLM extraction pipeline (batch size 5,
confidence threshold 0.65). Table~\ref{tab:corpus} reports the
preliminary extraction statistics.

\begin{table}[ht]
\centering
\caption{Corpus-scale graph construction statistics.}
\label{tab:corpus}
\small
\begin{tabular}{lr}
\toprule
\textbf{Metric} & \textbf{Value} \\
\midrule
Corpus passages           & 30,533 \\
Candidate pairs (cross-source, $k=8$) & 49,571 \\
LLM extraction batches    & 9,915 \\
Confidence threshold $\theta$ & 0.65 \\
Extraction model          & \texttt{deepseek-chat} \\
Estimated edges (at 40\% yield) & $\approx$19,800 \\
Graph file size (gzip JSON) & $<$10 MB \\
\bottomrule
\end{tabular}
\end{table}

Full extraction and edge-type distribution analysis are ongoing;
quantitative retrieval evaluation against HotpotQA and MuSiQue
baselines is reported in Section~\ref{sec:future}.

\subsection{Qualitative Analysis}
\label{sec:qualitative}

We examined a random sample of 50 extracted edges from the governance
corpus and classified them as correct, incorrect, or marginal by manual
annotation. Of these, 39 (78\%) were judged correct, 7 (14\%) marginal
(the relationship exists but the type assignment was debatable), and
4 (8\%) incorrect. The most common error type was conflating
\textsc{qualifies} with \textsc{contrasts\_with} when a passage adds
significant exceptions to another.

Edge type distribution in the extracted sample was:
\textsc{supports} 32\%, \textsc{specializes} 20\%,
\textsc{qualifies} 18\%, \textsc{derives\_from} 12\%,
\textsc{contrasts\_with} 7\%, \textsc{refutes} 5\%,
\textsc{implements} 3\%, \textsc{supersedes} 2\%,
\textsc{generalizes} 1\%, \textsc{exemplifies} $<$1\%.
This distribution is consistent with a corpus of technical standards,
where supportive elaboration dominates and outright contradiction is rare.

% ── 6. Planned Evaluation ─────────────────────────────────────────────────────

\section{Planned Evaluation}
\label{sec:future}

The following evaluations are in preparation:

\paragraph{Multi-hop QA benchmarks.}
We will evaluate PRISM against dense retrieval and HippoRAG
\citep{gutierrez2024hipporag} on HotpotQA \citep{yang2018hotpotqa},
MuSiQue \citep{trivedi2022musique}, and 2WikiMultiHopQA
\citep{ho2020constructing}. These benchmarks require connecting
information across documents, which is exactly the setting where
graph-augmented retrieval has demonstrated gains. Primary metrics:
Recall@$K$ and Answer F1.

\paragraph{Epistemic structure quality.}
We will construct a synthetic adversarial QA set with known ground-truth
epistemic structure: 200 questions over a corpus with deliberately planted
contradictions (requiring \textsc{contrasting} bucket), supersession
events (requiring \textsc{superseded} bucket), and multi-source
convergence (testing convergence scoring). This evaluates what existing
benchmarks cannot: whether PRISM correctly identifies the \emph{epistemic
role} of retrieved passages, not just their relevance.

\paragraph{LLM reasoning quality study.}
We will conduct a human evaluation comparing LLM answers generated from
(a) flat top-5 vector retrieval context and (b) PRISM structured epistemic
context, on questions for which the corpus contains conflicting information.
The hypothesis is that LLMs given PRISM's bucketed output will produce
fewer factual errors and more appropriately hedged answers on
contradiction-containing queries.

\paragraph{Baseline comparison.}
Full head-to-head comparison against HippoRAG, GraphRAG, and LightRAG
on HotpotQA is the primary benchmark target. We expect PRISM's advantage
to be largest on queries where (a) the corpus contains genuine epistemic
conflict or qualification, and (b) multi-document synthesis is required.

% ── 7. Limitations ────────────────────────────────────────────────────────────

\section{Limitations}
\label{sec:limitations}

\paragraph{Graph quality is LLM-dependent.}
The epistemic graph is produced by an LLM that can misclassify relationships.
A passage tagged \textsc{supports} when it only tangentially relates, or
\textsc{refutes} when it merely offers a different framing, degrades retrieval
quality. Graph quality is correlated with chunk quality, chunking strategy,
and LLM capability. We recommend reviewing extracted edges before
trusting them in high-stakes applications.

\paragraph{Confidence scores are not calibrated probabilities.}
The confidence values returned by the extraction LLM are self-reported
estimates, not calibrated. They are used as a filter and as edge weights,
but should not be treated as precise measures of relationship strength.

\paragraph{Noisy edges degrade retrieval.}
If the graph contains many false-positive edges, spreading activation
propagates to irrelevant nodes. In the worst case, PRISM can perform
worse than pure vector retrieval. The fallback to vector-only retrieval
for empty graphs does not help when the graph is large but noisy.

\paragraph{Cross-source assumption.}
PRISM is optimised for multi-source corpora. The \texttt{cross\_source\_only}
default reduces noise but means single-source corpora see fewer extracted
edges and less epistemic structure.

\paragraph{Graph construction cost.}
The one-time build requires thousands of LLM calls (hours for large corpora)
and has no incremental update path. Adding new documents currently
requires a full rebuild.

\paragraph{No QA evaluation benchmark yet.}
The retrieval quality improvements claimed in this paper are supported by
qualitative analysis and synthetic benchmarks only. Head-to-head comparison
on standard QA datasets is in progress (Section~\ref{sec:future}).

% ── 8. Conclusion ──────────────────────────────────────────────────────────────

\section{Conclusion}
\label{sec:conclusion}

We have presented PRISM, a retrieval framework that augments vector search
with a typed epistemic knowledge graph and spreading activation to produce
structured, epistemically annotated retrieval results. The core insight is
that knowledge retrieval should return not a ranked list of similar passages
but a structured account of how retrieved knowledge relates to the query:
what supports it, what challenges it, what qualifies it, what supersedes it.

PRISM operationalises this through three novel components: a ten-type
epistemic edge taxonomy with formal propagation semantics, a convergence-scoring
variant of spreading activation that rewards multi-perspective consensus, and an
epistemic bucketing output schema that structures retrieval results for LLM
consumption. The system requires no re-indexing of existing vector stores and
degrades gracefully to standard vector retrieval when no graph is available.

The framework is available as an open-source Python package at
\url{https://pypi.org/project/prism-rag/} and
\url{https://github.com/MadMando/prism}. We invite the community to
evaluate PRISM on their own corpora, contribute edge type extensions,
and collaborate on the quantitative benchmarks described in
Section~\ref{sec:future}.

% ── References ─────────────────────────────────────────────────────────────────

\bibliographystyle{plainnat}
\bibliography{references}

% Inline bibliography for arXiv (no separate .bib required)
\begin{thebibliography}{99}

\bibitem[Brachman \& Levesque, 1985]{brachman1985overview}
Brachman, R.~J., \& Levesque, H.~J. (1985).
\newblock Readings in knowledge representation.
\newblock Morgan Kaufmann.

\bibitem[Collins \& Loftus, 1975]{collins1975spreading}
Collins, A.~M., \& Loftus, E.~F. (1975).
\newblock A spreading-activation theory of semantic processing.
\newblock \emph{Psychological Review}, 82(6), 407--428.

\bibitem[Edge et al., 2024]{edge2024graphrag}
Edge, D., Trinh, H., Cheng, N., Bradley, J., Chao, A., Mody, A.,
  Truitt, S., \& Larson, J. (2024).
\newblock From local to global: A graph {RAG} approach to query-focused
  summarization.
\newblock \emph{arXiv:2404.16130}.

\bibitem[Fagin et al., 2003]{fagin2003reasoning}
Fagin, R., Halpern, J.~Y., Moses, Y., \& Vardi, M.~Y. (2003).
\newblock \emph{Reasoning About Knowledge}.
\newblock MIT Press.

\bibitem[Guo et al., 2024]{guo2024lightrag}
Guo, Z., Yao, L., Yang, M., Ji, J., Wang, P., Chen, X., \& Lin, W. (2024).
\newblock {LightRAG}: Simple and fast retrieval-augmented generation.
\newblock \emph{arXiv:2410.05779}.

\bibitem[Guti\'{e}rrez et al., 2024]{gutierrez2024hipporag}
Guti\'{e}rrez, B.~J., Shu, Y., Gu, Y., Yasunaga, M., \& Su, Y. (2024).
\newblock {HippoRAG}: Neurobiologically inspired long-term memory for large
  language models.
\newblock \emph{Advances in Neural Information Processing Systems (NeurIPS)}.
\newblock \emph{arXiv:2405.14831}.

\bibitem[Ho et al., 2020]{ho2020constructing}
Ho, X., Nguyen, A.-K.~D., Sugawara, S., \& Aizawa, A. (2020).
\newblock Constructing a multi-hop {QA} dataset for comprehensive evaluation
  of reasoning steps.
\newblock \emph{Proceedings of COLING 2020}.

\bibitem[Karpukhin et al., 2020]{karpukhin2020dpr}
Karpukhin, V., O\u{g}uz, B., Min, S., Lewis, P., Wu, L., Edunov, S.,
  Chen, D., \& Yih, W.-t. (2020).
\newblock Dense passage retrieval for open-domain question answering.
\newblock \emph{Proceedings of EMNLP 2020}.

\bibitem[Lewis et al., 2020]{lewis2020rag}
Lewis, P., Perez, E., Piktus, A., Petroni, F., Karpukhin, V., Goyal, N.,
  K\"{u}ttler, H., Lewis, M., Yih, W.-t., Rockt\"{a}schel, T., Riedel, S.,
  \& Kiela, D. (2020).
\newblock Retrieval-augmented generation for knowledge-intensive {NLP} tasks.
\newblock \emph{Advances in Neural Information Processing Systems (NeurIPS)}.

\bibitem[McClelland et al., 1995]{mcclelland1995complementary}
McClelland, J.~L., McNaughton, B.~L., \& O'Reilly, R.~C. (1995).
\newblock Why there are complementary learning systems in the hippocampus
  and neocortex: Insights from the successes and failures of connectionist
  models of learning and memory.
\newblock \emph{Psychological Review}, 102(3), 419--457.

\bibitem[Pollock, 1995]{pollock1995cognitive}
Pollock, J.~L. (1995).
\newblock \emph{Cognitive Carpentry: A Blueprint for How to Build a Person}.
\newblock MIT Press.

\bibitem[Teyler \& DiScenna, 1986]{teyler1986hippocampal}
Teyler, T.~J., \& DiScenna, P. (1986).
\newblock The hippocampal memory indexing theory.
\newblock \emph{Behavioral Neuroscience}, 100(2), 147--154.

\bibitem[Trivedi et al., 2022]{trivedi2022ircot}
Trivedi, H., Balasubramanian, N., Khot, T., \& Sabharwal, A. (2022).
\newblock Interleaving retrieval with chain-of-thought reasoning for
  knowledge-intensive multi-step questions.
\newblock \emph{Proceedings of ACL 2023}.

\bibitem[Trivedi et al., 2022]{trivedi2022musique}
Trivedi, H., Balasubramanian, N., Khot, T., \& Sabharwal, A. (2022).
\newblock {MuSiQue}: Multihop questions via single-hop question composition.
\newblock \emph{Transactions of the ACL}.

\bibitem[Yang et al., 2018]{yang2018hotpotqa}
Yang, Z., Qi, P., Zhang, S., Bengio, Y., Cohen, W.~W., Salakhutdinov, R.,
  \& Manning, C.~D. (2018).
\newblock {HotpotQA}: A dataset for diverse, explainable multi-hop
  question answering.
\newblock \emph{Proceedings of EMNLP 2018}.

\end{thebibliography}

\end{document}
